\chapter{15}

\section{Dienstag, 21. August}



Louis ahnte Jess’ Blick im Rücken, als er sich in der Dunkelheit des
frühen Morgens geräuschlos davonschlich.
Gut sichtbar hatte er für sie einen 20-Euroschein bei seinem Schlafplatz
zurückgelassen. Ohne Notiz. Sie würde das Geld gebrauchen können. Eine
billige Geste. Etwas Besseres war ihm nicht eingefallen. Immerhin, hatte
er sich eingeredet. Louis hoffte, dass sie klug genug war, sich etwas zu
essen zu kaufen.

\sectionbreak

Die Sonne brannte ohne Gnade zwischen den Hausfassaden auf den Asphalt hinunter. 
Bald würde die Stadt erneut zum Hitzekessel werden. Louis schleppte sich ziellos durch die Strassen. 
Bis er schlagartig stehen blieb und sich an den Kopf fasste. Ein Passbild, er musste bis zum Abend zu einem neuen Passbild kommen. 
Er hatte Glück. Eine ältere Frau wechselte ihm einige Euros zu Schweizer Franken und wünschte ihm einen schönen Tag.
Die Fotobox befand sich in unmittelbarer Nähe.
Der an der Aussenwand des Automaten angebrachte Spiegel war nahezu
blind. Louis kämmte sein feuchtes Haar, putzte die Brillengläser und
setzte das Gestell auf.
Umständlich zwängte er sich in die kleine Box, warf die Münzen ein und
wartete auf den ersten Blitz. Er machte ein ernstes Gesicht. Zwei Bilder
würden geschossen, las er auf der Anleitung unterhalb des Bildschirms.
Dann, drei Minuten warten. Ungeduldig stellte er sich vor den Automaten und trank seine Flasche leer.

Schliesslich spuckte das Gerät ein rechteckiges Stück Fotopapier mit
zwei mal zwei Bildern aus. Er wartete die Lufttrocknung ab. Dann
fingerte er die Aufnahmen sorgfältig aus dem kleinen Kästchen und nahm
seine Brille ab. Er schaute einem Fremden ins Gesicht. Verdattert
streckte er die Bilder auf Armeslänge von sich. Er sah schrill aus. Doch
als Passbilder schienen die Aufnahmen in Ordnung.
Er holte das Sackmesser aus dem Rucksack, schnitt die Fotos mit der
kleinen Schere auseinander und steckte sie ein. Erwartungsvoll sehnte er die Nacht herbei. 

Während er sich Richtung See aufmachte, fragte er sich, ob er Mascha finden würde. Wer der Typ war. Wie er ihm begegnen würde. Er hoffte, nirgendwo angehalten oder aufgegriffen zu werden, und überhaupt, dass er gar nicht erst gesucht wurde. Doch wenn, was er tun würde. Schon nach den ersten Schritten wusste Louis, wenigstens seine zweite Mission «Zeit totschlagen» würde leichter gehen. 
Im  Vergleich zum Vortag kannte er sich in den Strassen bereits aus. Zürich war so viel kleiner als Mailand. 
Die Hitze trieb ihn auf direktem Weg auf eine weitläufige Wiese an der Seepromenade. Am Ende des Weges und einige Meter vom Ufer entfernt entdeckte er ein kleines Kaffee. Es war voll besetzt. Louis kaufte sich einen Kaffee im Becher, schlenderte damit über die Wiese und richtete sich schliesslich unter dem Schatten eines Baumes ein. Er war verblüfft, wie selbstverständlich es sich anfühlte, zwischen den unzähligen Menschen zu sitzen. Einzeln und in Gruppen schienen sie mit sich selbst beschäftig zu sein und der Kaffee schmeckte. Hier liess man sich offenbar in Ruhe und er sah keinen Anlass, sich zu verstecken. Umgeben von Stimmen, Gelächter, Kindergeschrei und einzelnen Musikfetzen legte sich Louis hin, schloss die Augen und schlummerte ein. Sein letzter Gedanke hatte Jess gegolten.

\sectionbreak

Endlich war es kurz nach dreiundzwanzig Uhr und dunkel.
Louis machte sich auf den Weg an die Langstrasse. Er würde die Adressen
abklappern, die ihm Jess genannt hatte. Was, wenn er Mascha nicht fand?
Oder ihn nicht erkannte?
Er suchte nach den genannten Strassenschildern, konnte sich leicht
orientieren und staunte, wie präzise Jess den Weg beschrieben hatte. 
Nach einer Stunde war er verzweifelt. Die Hoffnung schwand. Immerhin, 
er war einer Polizeikontrolle aus dem Weg gegangen und hatte sich mehrmals vor patrouillierenden
Polizeiwagen in Deckung gebracht.

Doch dann - ein Mann stieg aus einem weissen Mercedes, direkt vor der «Gay-Bar», die
ihm Jess genannt hatte. Hier stand er bereits das zweite Mal in dieser
Nacht.
Erst konnte er ihn nicht genau erkennen. Der Mann beugte sich zum
Fenster des Fahrers hinunter und schäkerte mit ihm. Louis sah, dass er etwas entgegennahm 
und dem Fahrer anschliessend lasziv nachwinkte.
Dann drehte er sich um, und Louis war sich ziemlich sicher. Das musste
Mascha sein.
Vor der Bar standen einige Männer herum. Ein Jüngerer schaute Louis
interessiert und direkt an. Louis blickte angestrengt weg.
Der Mann schritt auf den Eingang der Bar zu. Das war der Moment. Jetzt
musste Louis ihn ansprechen. Das rot-gelb blinkende Reklameschild im
Hintergrund konnte er nur schwer ausblenden.

«Hey, Mascha?»

Louis blieb vor ihm stehen. Dessen Gesichtsausdruck wechselte in
Sekundenschnelle von gedankenabwesend zu zuckersüss.

«Ja, bin ich,» flötete er.

Louis schaute Mascha fasziniert an. Kleidung und Gesten schienen auf
einander abgestimmt. Sein Blick war geradezu betörend. Hemmungslos.

«Darf ich dich was fragen?» sagte Louis.

Mascha zauberte sich ein Lächeln auf sein Gesicht. Er trat etwas näher
an Louis heran. Seine Stimme klang sonor und warm.

«Aber sicher, mein Süsser. Und keine Angst. Geht leicht. Bist neu hier.
Das erste Mal?»

Er fixierte Louis mit einem Blick, als wollte er ihn hypnotisieren.
Louis fühlte Panik aufkommen. Irgend etwas sträubte sich in ihm. Doch es
musste sein. Sofort.

«Ich komme von Jess, vielleicht kannst du mir helfen?»

Wieder veränderte sich Maschas Ausdruck in einem Sekundenbruchteil. 
Louis glaubte, einen Anflug von Besorgnis zu erkennen. Jede Form von Anmache war weg.
Mascha schaute ihn ungeduldig an.

«Jess? Was ist mit ihr?»

«Nichts, alles ok,» Louis versuchte locker zu klingen, «ich bin ein Freund und sie sagte mir, vielleicht kannst du mir
helfen?»

Er wischte sich den Schweiss von der Stirn.
Maschas Blick war nun forschend. Er kniff seine Augen zusammen und
wirkte vorsichtig. Mit einer leichten Handbewegung wies er Louis in eine
kleine Seitengasse. Er schaute nervös und gleichzeitig dezent in alle
Richtungen. Hier war keiner.

«Code?» fragte Mascha.

«Süssholz.»

«Gut.»

Maschas Augen blitzten. Er entspannte sich sichtlich, nickte und war
wieder zugänglicher.
Aus einem geöffneten Fenster klang Musik.
Jemand leierte am Klavier Beethovens \emph{«Für Elise»} herunter.

\qrblock{Beethoven \emph{«Für Elise»}}{media/QR-Codes/015_Für_Elise_Beethoven_Spotify.png}

Louis wurde übel. Sein Kopf schmerzte. Und wie er hier vor Mascha stand,
begriff er, dass er es endlich verstehen musste. Er stand direkt davor,
sich in kriminelle Handlungen zu verstricken. Oder war eigentlich
bereits mittendrin. Er dachte an Giorgia. Nur kurz, denn urplötzlich
krachten schwere Regentropfen auf sie hinunter. Sie fühlten sich spitz
an. Louis atmete das bekannte Gemisch von trockenem Teer und frischem
Regen ein. Es schien, als giesse sich der Himmel leer. Louis hielt seine
Hände schützend über die Augen. Mascha packte ihn plötzlich am Arm und
zog ihn unter das schmale Dach eines Hauseingangs. Nach kürzester Zeit
waren der Platz und all die Menschen rund um sie herum platschnass.
Louis holte tief Luft.

«Ich brauche einen neuen Pass. Kannst du mir einen besorgen?»

Es war raus. Mascha zögerte. Mit Seitenblicken nach links und rechts
wischte er sich die Wassertropfen aus dem Gesicht.

«Ja, kann ich. Passbild dabei?»

Er artikulierte sich überdeutlich, gepresst leise, zwischen das Prasseln
des Regens und hinein in das gedämpfte Stimmengewirr um sie herum.
Wieder schaute er konzentriert in die Umgebung, dann wieder zu Louis.
Sie standen so nahe beisammen, dass Louis Maschas schweren Duft eines
Herrenparfums roch. Die Intensität nahm ihm den Atem. Er schnappte
flüchtig nach Luft.

«Ja, hab ich hier.»

Dabei zog er das Bild aufgeregt aus der Hosentasche und reichte es ihm.
Mascha verglich Louis’ Gesicht mit dem Passbild.

«Zwei Tage. Übermorgen hier, gleiche Zeit,» seine knappen Worte hörten
sich an wie dumpfe Schüsse, er blickte auf sein Handy, «fünfzehn Minuten nach Mitternacht.»

Er musterte Louis von oben nach unten.

«Schweizer, ca. 30 Jahre, komplett nass,» sein Lacher kam unerwartet,
«passt, oder? Mindestens ein Elternteil könnte aus Italien, der Türkei
oder aus Spanien sein?"

Louis nickte mehrmals. Er konnte kaum glauben, wie geübt Mascha das
Geschäft abwickelte.

«Ein spanischer Name wäre gut.»

Er schaute Mascha unruhig an.

«Ist kein Wunschkonzert, mein Lieber. Ich werde fragen, was da ist. Und
machbar. Ich hoffe, ich find was Passendes.»

Nun veränderten sich Haltung und Gesichtsausdruck schon wieder. Das
charmante Lächeln war weg.
Er schaute Louis diesmal noch eindringlicher und geradezu herrisch an.

«Und kein Wort zu niemandem, niemals, verstehst du mich? Sonst lernst du
meine andere Seite kennen,» Mascha kam Louis noch näher, sodass er
dessen Atem in seinem Gesicht fühlen konnte, «ich kann sehr, sehr böse
werden.»

Er ballte die Faust. Louis’ Blut pulsierte spürbar in seinen Adern. Der
Regen hatte nachgelassen.

«Klar.»

Louis’ bebte.

«Also, Deal. Siebenhundertfünfzig, weil du’s bist, in bar bei der
Übergabe.»

Mascha trat einen Schritt zurück. Das war viel Geld. Louis unterstand
sich, den Preis in Frage zu stellen.

«Ich habe Euro.»

Maschas Wesen mutierte in Windeseile zurück zum freundlichen Typen.

«Passt.»

Der Mann war der geborene Schauspieler.
Louis wollte sich gerade zum Gehen umdrehen, als er ihn an der Schulter
zurückhielt.
